\documentclass[10pt,a4paper]{amsart}
\usepackage[margin=19mm]{geometry}
\usepackage{amsmath,amssymb,mathtools}
\usepackage[scr=rsfs,cal=euler]{mathalfa}
\usepackage{xfrac}
\usepackage[unicode,hidelinks,bookmarks=false]{hyperref}
\newcommand{\ie}{i.e.\ }
\newcommand{\cf}{cf.\ }
\newcommand{\resp}{respectively\ }
\newcommand{\Subsection}{\subsection}
\providecommand{\nobreakdash}{\nobreak\hbox{-}}
\setlength{\emergencystretch}{2em}
\multlinegap0pt
\usepackage{bm}
\usepackage{bbm}
\usepackage{dsfont}
\usepackage{xspace}
\usepackage{cancel}
\usepackage{mathtools}
\usepackage{amsthm}
\makeatletter
\@ifundefined{lemma}{\newtheorem{lemma}{Lemma}}{}
\makeatother
\newcommand{\R}{\mathbb R}
\newcommand{\Z}{\mathbb Z}
\title[A discrete approach to the functional Zhang's projection ine]{A discrete approach to the functional Zhang's projection inequality}
\author{David Alonso-Guti\'errez, Julia S\'anchez-Loscertales}
\date{Source: arXiv:2607.20018v1 (CC BY 4.0)}
\begin{document}
\maketitle

\noindent\emph{Source and licence.} A discrete approach to the functional Zhang's projection inequality. Version arXiv:2607.20018v1 (CC BY 4.0), distributed under the Creative Commons Attribution 4.0 International licence, as recorded in the arXiv licence metadata for this exact version and shown on the abstract page https://arxiv.org/abs/2607.20018v1. Full rights notice: licenses/arxiv-2607.20018-CC-BY-4.0.txt.\par\smallskip

\noindent\textit{Editorial scope.} Counted unit: the Lemma whose source label is \texttt{lem:LogConcaveIntegralDiscreteToContinuous} in the article (source0.tex lines 447--456, source label \texttt{lem:LogConcaveIntegralDiscreteToContinuous}), delivered together with the article's own complete original proof of it (source0.tex lines 458--499, one \texttt{proof} environment of 42 lines). No mathematics is written, repaired or re-derived: statement and proof are the article's own text line for line. Required context retained uncounted: eq:LatticePointAndVolume (lines 288--290); discrete\_continuous\_volume (lines 293--295); eq:nu(L)Integral (lines 374--376); eq:barnu(L)Integral (lines 378--380); eq:InclusionSuperlevelSetsAsplund (lines 405--407); eq:InverseInclusionSuperLevelSetsAsplund (lines 414--416). Every definition, display and label of those lines is retained unchanged. Omitted from this package and not redistributed: 1--287, 291--292, 296--373, 377--377, 381--404, 408--413, 417--446, 457--457, 500--1860. Nothing inside a retained excerpt is omitted or altered: every retained body is delivered exactly as the article's lines are written, so no omission receipt of any kind accompanies this package. Citations: the retained excerpts carry no citation command at all, so no bibliography entry of the article is transcribed and no reference is redistributed. Reference adaptation: none was needed and none was made; every reference inside the delivered unit resolves inside it, so no unresolved reference or citation is produced. Printed slips kept literal: no printed word, symbol, hypothesis or number of the counted statement or of its proof was altered. Layout and class: the article's own class is \texttt{amsart} with packages including amsfonts, color, latexsym, amssymb, amsmath, amsthm, epsfig, graphicx, colortbl, url, hyperref, mathtools, changes, mathptmx; the wrapper is the lane's pinned \texttt{amsart} editorial wrapper at 10pt on A4, which restates the theorem-like environments the retained text needs and adds the article's own macro definitions that the retained text needs (\texttt{\textbackslash R}, \texttt{\textbackslash Z}), with extra wrapper packages (bm, bbm, dsfont, xspace, cancel, mathtools). The delivered numbering is the unit's own. Assets: no retained span contains a figure environment, a graphics-inclusion command, a PDF-page inclusion, a table or a TikZ picture; no image file is copied into this package, the package declares no asset dependency and no image file is redistributed with it. Licence: version arXiv:2607.20018v1 of this article is distributed under the Creative Commons Attribution 4.0 International licence, as recorded in the arXiv licence metadata for this exact version and shown on the abstract page https://arxiv.org/abs/2607.20018v1. A full rights notice is delivered at licenses/arxiv-2607.20018-CC-BY-4.0.txt. Characters: every retained excerpt is pure ASCII, so the delivered engine, XeLaTeX with its default Latin Modern text font, reports no missing character.
\begin{equation}\label{eq:LatticePointAndVolume}
G_n(A)=\left|(A\cap\Z^n)+\frac{1}{2}B_\infty^n\right|\leq\left|A+\frac{1}{2}B_\infty^n\right|.
\end{equation}

\begin{equation}\label{discrete_continuous_volume}
\lim_{\lambda \to \infty} \frac{G_n(\lambda K+M)}{\lambda^n} = |K|_n.
\end{equation}

\begin{equation}\label{eq:nu(L)Integral}
\nu(L(f))=\int_0^\infty e^{-t} G_n(L_t(f))dt=\int_0^1\int_{\R^n}\chi_{\left\{x\in\R^n\,:\,\frac{f(x)}{\Vert f\Vert_\infty} \geq s \right\}}(x)dG_n(x)ds=\int_{\R^n}\frac{f(x)}{\Vert f\Vert_\infty}dG_n(x),
\end{equation}

\begin{equation}\label{eq:barnu(L)Integral}
\overline{\nu}(L(f))=\int_0^\infty e^{-t} |L_t(f)|dt=\int_0^1\int_{\R^n}\chi_{\left\{x\in\R^n\,:\,\frac{f(x)}{\Vert f\Vert_\infty}\geq s\right\}}(x)dxds=\int_{\R^n}\frac{f(x)}{\Vert f\Vert_\infty}dx,
\end{equation}

\begin{equation}\label{eq:InclusionSuperlevelSetsAsplund}
L_{t}(f)+C\subseteq \bigcap_{\varepsilon>0}L_{t+\varepsilon}(f\star \chi_C)=L_t(f\star\chi_C).
\end{equation}

\begin{equation}\label{eq:InverseInclusionSuperLevelSetsAsplund}
L_t(f\star\chi_C)\subseteq\bigcap_{\varepsilon>0}\overline{L_{t+\varepsilon}(f)}+\overline{C}.
\end{equation}

\begin{lemma}\label{lem:LogConcaveIntegralDiscreteToContinuous}
Let $f\in\mathcal{F}(\R^n)$ be continuous on its support and, for any $\lambda>0$, let  $f_\lambda\in\mathcal{F}(\R^n)$ be the function defined as $f_\lambda(x)= f\left(\frac{x}{\lambda}\right)$. Let also $M$ be a bounded set containing the origin. Then,
$$
\lim_{\lambda\to\infty}\frac{1}{\lambda^n}\int_{\R^n}f_\lambda\star\chi_{M}(x)dG_n(x)=\int_{\R^n}f(x)dx.
$$
In particular, taking $M=\{0\}$,
$$
\lim_{\lambda\to\infty}\frac{1}{\lambda^n}\int_{\R^n}f_\lambda(x)dG_n(x)=\int_{\R^n}f(x)dx.
$$
\end{lemma}

\begin{proof}
Notice that, since $\Vert f_\lambda\star\chi_M\Vert_\infty=\Vert f_\lambda\Vert_\infty=\Vert f\Vert_\infty$, by \eqref{eq:nu(L)Integral},
$$
\lim_{\lambda \to \infty} \int_{\R^n} \frac{f_\lambda\star\chi_M (x)}{\Vert f\Vert_\infty} dG_n(x)=\lim_{\lambda\to\infty}\int_0^\infty e^{-t} \frac{G_n (L_t (f_\lambda\star\chi_M))}{\lambda^n}dt.
$$

On the one hand, since $f$ is continuous on its support, we have that $L(f)$ is closed. Thus, by \eqref{eq:InclusionSuperlevelSetsAsplund} and \eqref{eq:InverseInclusionSuperLevelSetsAsplund}, we have that for every $t\geq0$ and every $\lambda>0$
$$
L_t(f_\lambda)+M\subseteq L_t(f_\lambda\star\chi_M)\subseteq\bigcap_{\varepsilon>0}\overline{L_{t+\varepsilon}(f_\lambda)}+\overline{M}=\bigcap_{\varepsilon>0}L_{t+\varepsilon}(f_\lambda)+\overline{M}=L_t(f_\lambda)+\overline{M}.
$$
Therefore, for every $t\geq0$ and every $\lambda>0$,
$$
\lambda L_t(f)+M\subseteq L_t(f_\lambda\star\chi_M)\subseteq \lambda L_t(f)+\overline{M}.
$$
Since by \eqref{discrete_continuous_volume}, for every $t\geq0$ we have that
$$
\lim_{\lambda\to\infty}e^{-t} \frac{G_n (\lambda L_t (f)+M)}{\lambda^n}=e^{-t} |L_t (f)|\quad\textrm{and}\quad\lim_{\lambda\to\infty}e^{-t} \frac{G_n (\lambda L_t (f)+\overline{M})}{\lambda^n}=e^{-t} |L_t (f)|,
$$
we obtain that for every $t\geq0$
$$
\lim_{\lambda\to\infty}e^{-t}\frac{G_n (L_t(f_\lambda\star\chi_M))}{\lambda^n}=e^{-t}|L_t(f)|.
$$


On the other hand, since $M$ is bounded, there exists $R>0$ such that $\overline{M}\subseteq RB_\infty^n$. Therefore,  using \eqref{eq:LatticePointAndVolume} and \eqref{eq:InclusionSuperlevelSetsAsplund}, for every $t>0$ and every $\lambda\geq R+\frac{1}{2}$
\begin{eqnarray*}
G_n(L_t(f_\lambda\star\chi_M))&\leq& G_n(\lambda L_t(f)+\overline{M})\leq\left|\lambda L_t(f)+\overline{M}+\frac{1}{2}B_\infty^n\right|\cr
&\leq&\left|\lambda L_t(f)+\left(R+\frac{1}{2}\right)B_\infty^n\right|\leq\left|\lambda L_t(f)+\lambda B_\infty^n\right|=\lambda^n|L_t(f)+B_\infty^n|\cr
&\leq&\lambda^n L_t(f\star\chi_{B_\infty^n}).
\end{eqnarray*}
 Thus, for every $t\geq0$ and every $\lambda>R+\frac{1}{2}$ we have
$$
e^{-t} \frac{G_n(L_t(f_\lambda\star\chi_M))}{\lambda^n}\leq e^{-t}\left|L_t (f\star \chi_{B_\infty^n})\right|,
$$
which, taking into account \eqref{eq:barnu(L)Integral}, is integrable on $[0,\infty)$, since $f\star \chi_{B_\infty^n}\in\mathcal{F}(\R^n)$.

By the dominated convergence theorem,
$$
\lim_{\lambda\to\infty}\int_0^\infty e^{-t} \frac{G_n(L_t(f_\lambda\star\chi_M))}{\lambda^n}dt=\int_0^\infty e^{-t} |L_t (f)|dt=\int_{\R^n}\frac{f(x)}{\Vert f\Vert_\infty}dx,
$$
which proves the lemma.
\end{proof}

\end{document}
